\chapter*{Résumé}
\addcontentsline{toc}{chapter}{Résumé}

Cette alternance de deux années au sein du service Recherche \& Développement de Souchier-Boullet s'inscrit dans un contexte industriel où la fiabilité des données techniques, la maîtrise des exigences réglementaires et la réactivité opérationnelle sont étroitement liées. L'entreprise conçoit et industrialise des solutions de compartimentage coupe-feu, de désenfumage naturel et, pour certaines gammes, des produits intégrant également des performances acoustiques. Dans ce cadre, mes missions se sont organisées autour de deux axes complémentaires : la maintenance et l'amélioration de modèles 3D sous Autodesk Inventor, d'une part, et le développement de configurateurs produit reposant sur une architecture SQL et une interface Infor Design Studio, d'autre part.

Le premier volet de mon travail a consisté à traiter des fiches de modification issues de la production et du bureau d'études, afin d'assurer la cohérence entre les modèles numériques, les plans techniques et la réalité de fabrication. Le second volet, plus structurant à l'échelle de l'entreprise, a porté sur la transformation d'une logique de configuration historiquement fondée sur la lecture manuelle de procès-verbaux feu et de tableaux Excel en un système relationnel robuste, capable de formaliser les règles métier, d'automatiser des calculs et de sécuriser la sélection des variantes compatibles.

La méthodologie construite en A4 sur la gamme Coveraccess a été consolidée par une première extension à DOORSLIDE, puis approfondie en A5 avec DOORSTEEL, PHONI LUX-AIR-PASS, les portes ONYX et MULTIDOOR, et les rideaux ScreenPass et CURTEX V2. Cette progression a confirmé la pertinence d'une architecture générique, tout en mettant en évidence la diversité des contraintes produit, qu'elles soient géométriques, réglementaires, mécaniques ou acoustiques. Une avancée marquante de la seconde année a consisté à relier dynamiquement le configurateur SQL à Autodesk Inventor, afin d'ouvrir automatiquement le modèle 3D correspondant aux paramètres choisis dans l'interface.

Le mémoire montre que ces développements ont produit des effets concrets : meilleure capitalisation des connaissances techniques, réduction des risques d'erreur liés aux manipulations manuelles, amélioration du dialogue entre les équipes et préparation d'une continuité numérique plus fiable. Il met aussi en lumière les limites d'une intégration reposant sur des conventions implicites entre outils. Le chapitre scientifique approfondit précisément cette question en proposant un contrat d'interface formel entre le configurateur SQL et le modeleur CAO, puis en évaluant un mécanisme de validation capable de détecter les ruptures de compatibilité avant exécution.

\paragraph{Mots-clés :} SQL, configurateur produit, CAO, Autodesk Inventor, compartimentage, désenfumage, interopérabilité, Infor Design Studio, transformation numérique industrielle.
